\chapter{Errors, Exceptions, and Unwinding across Language Boundaries}
\label{ch:ffi-errors}

A callback can fail after native code has acquired resources and before it publishes its result.
If the runtime jumps past the native cleanup path, those resources may leak or remain locked.
If the adapter converts every failure into a generic message, it can lose the original exception and its diagnostic context.
The boundary must translate both the failure representation and the path by which control returns.

\section{Choose an explicit error channel}
\label{sec:ch74-channels}
The case study uses a status return plus output parameters.
Success initializes the output; failure leaves it untouched.
This separates the full numeric result domain from failure indicators and permits a callback to return through ordinary C frames.
Table~\ref{tab:ch74-channels} compares common channels and their obligations.

\begin{table}[htbp]
\caption{Each error channel imposes a different interpretation and lifetime rule.}
\label{tab:ch74-channels}
\centering
\begin{tabular}{ll}
\toprule
Channel & Boundary obligation \\
\midrule
status plus output & inspect status before reading output \\
sentinel result & reserve sentinel and identify auxiliary error state \\
thread-local error & capture promptly before another operation changes it \\
exception object & preserve ownership and the original failure identity \\
nonlocal transfer & establish a compatible catch/protection boundary \\
\bottomrule
\end{tabular}
\end{table}

A negative return is not universally an error; many native APIs legitimately return negative values.
Likewise, a nonzero thread-local error value after a successful call need not indicate that this call failed.
The adapter follows the particular API's success predicate, then reads any associated error state at the specified time.
Logging or cleanup can itself call library functions and overwrite that state, so capture the needed information first.

\section{Unwinding metadata is not a language contract}
\label{sec:ch74-unwind}
A platform unwinder can traverse frames using unwind metadata, but language-specific personality routines decide which handlers match and which cleanups run.
The Itanium C++ exception ABI separates search and cleanup phases and defines interactions between unwinding and language runtimes~\cite{itanium-2026-exceptions}.
Having compatible frame-walking information does not automatically make a C++ exception meaningful to Lua, Python, or an arbitrary JIT frame.

A practical C++ library boundary often catches exceptions within a C-callable wrapper, translates them to a status or owned error record, and returns normally.
The wrapper must cover every throwing operation, including allocation during error translation, and define a fallback for failures while constructing diagnostics.
An \code{extern "C"} declaration selects language linkage; it does not catch exceptions or prevent them from escaping.

Lua errors normally use nonlocal control transfer to an installed protected boundary~\cite{ierusalimschy-2026-lua}.
Jumping over C code skips ordinary local cleanup statements.
In the callback path, the adapter therefore calls the Lua function with \code{lua\_pcall}, converts failure into a C status, and re-raises the retained error object only after the native visitor returns.
Figure~\ref{fig:ch74-boundary} shows where the two control regimes meet.

\begin{figure}[htbp]
\centering
\begin{tikzpicture}[node distance=12mm]
\node[basic block] (failure) {managed callback raises};
\node[basic block,below=of failure] (catch) {protected adapter retains error};
\node[basic block,below=of catch] (return) {native frames return status};
\node[basic block,below=of return] (raise) {outer wrapper restores language failure};
\draw[cfg edge] (failure)--(catch); \draw[cfg edge] (catch)--(return);
\draw[cfg edge] (return)--(raise);
\end{tikzpicture}
\caption{The protected callback converts a nonlocal language failure into an ordinary native return.}
\label{fig:ch74-boundary}
\end{figure}

\section{Preserve the original callback failure}
\label{sec:ch74-preserve}
CPython reports failure from an object-returning API with a null result and a set exception indicator.
The adapter must propagate that failure rather than returning an apparently successful object while the indicator remains set~\cite{python-2026-capi}.
Our callback returns a native failure status without clearing the indicator; the outer wrapper then returns null.
The original exception object survives, and the test verifies its identity rather than only its message.

Lua permits error objects other than strings.
Our protected call leaves the original object on the stack, and the wrapper later invokes \code{lua\_error} with that object still at the top.
The test uses a table as the error marker and checks that the same table reaches the caller.
Replacing it with formatted text would lose this identity.

Algorithm~\ref{alg:ch74-propagate} describes the common structure.
It assumes cleanup does not replace the saved failure; a runtime that can invoke user finalizers during cleanup may need an explicit save-and-restore facility.
The adapter bookkeeping is $O(1)$ per callback plus runtime invocation, conversion, and cleanup costs.

\begin{algorithm}[htbp]
\caption{Return through native frames before propagating a managed callback error.}
\label{alg:ch74-propagate}
\begin{pseudocode}
\Require rooted callable, protected runtime entry, native status convention
\Ensure original managed failure is available to the outer wrapper
\State invoke the callable within the runtime's protected error boundary
\If{invocation or result conversion fails}
  \State retain the runtime's original error according to its ownership rules
  \State return native callback-failure status
\EndIf
\State copy the converted callback value to its output; return success
\Statex \textit{At the outer wrapper after native return:}
\State perform nonthrowing cleanup; propagate a retained failure if present
\end{pseudocode}
\end{algorithm}

Table~\ref{tab:ch74-trace} traces a failure at index two with output initially $-99$.
The local partial sum is not the public result.
This matters when a caller retries or distinguishes an unsuccessful operation from a completed operation returning a small value.
The callback's own external side effects, however, are not rolled back merely because the sum is withheld.

\begin{table}[htbp]
\caption{Failure preserves output but does not undo already executed callback effects.}
\label{tab:ch74-trace}
\centering
\begin{tabular}{lll}
\toprule
Event & Native state & Managed observation \\
\midrule
callback zero succeeds & local sum updated & first callback side effects persist \\
callback one succeeds & local sum updated & second callback side effects persist \\
callback two raises & failure returned & original error retained \\
visitor exits & output still $-99$ & no later callbacks \\
wrapper exits & native frame finished & original error propagated \\
\bottomrule
\end{tabular}
\end{table}

\section{Acquire resources transactionally}
\label{sec:ch74-cleanup}
Suppose a wrapper first exports a managed buffer, then allocates a native owner, then allocates the payload, and finally publishes a capsule.
Each failure point has a different cleanup suffix.
Algorithm~\ref{alg:ch74-acquire} records successful acquisitions and releases them in reverse order until publication transfers responsibility.
For $R$ resources this costs $O(R)$ bookkeeping and cleanup calls; actual destructors may cost more.

\begin{algorithm}[htbp]
\caption{Acquire and publish an owned result with complete failure cleanup.}
\label{alg:ch74-acquire}
\begin{pseudocode}
\Require ordered acquisitions with matching nonthrowing releases
\Ensure published owner or no outstanding acquired resources
\State initialize cleanup list to empty
\ForAll{required resources}
  \State attempt acquisition
  \If{acquisition fails}
    \State release recorded resources in reverse order; propagate failure
  \EndIf
  \State record its release action
\EndFor
\State publish owner; transfer its resources out of the cleanup list
\State release remaining temporary resources
\end{pseudocode}
\end{algorithm}

Figure~\ref{fig:ch74-cleanup} shows three acquisitions and their rollback dependencies.
This is a conceptual cleanup stack; the small C wrapper implements its few cases with explicit branches rather than a general runtime stack.
For the Lua wrapper, installing a finalizable userdata before native allocation supplies the persistent ownership path even if a later Lua operation raises.

\begin{figure}[htbp]
\centering
\begin{tikzpicture}[node distance=14mm]
\node[basic block] (a) {export input};
\node[basic block,below=of a] (b) {allocate owner};
\node[basic block,below=of b] (c) {allocate payload};
\draw[cfg edge] (a)--(b); \draw[cfg edge] (b)--(c);
\draw[cfg edge,dashed] (c.east) to[bend right=45] node[right,font=\small]{rollback}(b.east);
\draw[cfg edge,dashed] (b.west) to[bend left=45] node[left,font=\small]{rollback}(a.west);
\end{tikzpicture}
\caption{Failure releases only successfully acquired resources, in reverse dependency order.}
\label{fig:ch74-cleanup}
\end{figure}

\section{Cancellation is a protocol, not immediate destruction}
\label{sec:ch74-cancel}
Requesting cancellation does not prove a worker has stopped reading its buffer or that a queued callback cannot run.
A cancellable API needs a terminal acknowledgement after which ownership can be reclaimed.
Completion and cancellation can race, so exactly one transition should publish the terminal outcome and release each retained resource.
Cancellation of a wait also need not cancel the underlying operation.

The case study is synchronous and does not implement asynchronous cancellation.
Its callback failure simply stops further loop iterations and returns.
The retirement protocol in Chapter~\ref{ch:ffi-callbacks} supplies the additional quiescence requirement that an asynchronous variant would need before freeing registration state.

\section{Exercises}
\label{sec:ch74-exercises}
\begin{enumerate}
\item\label{ex:ch74-status} Why must a caller avoid reading the result after a range status when the output initially contains $99$?
\item\label{ex:ch74-identity} A Lua callback raises a table object.
Explain the observable difference between re-raising that object and raising its string representation.
\item\label{ex:ch74-cleanup} Enumerate cleanup after failure at each of owner allocation, payload allocation, and capsule publication, assuming an input export was acquired first.
\item\label{ex:ch74-cancel} A worker still holds a buffer when cancellation is requested.
Name the event required before freeing it and explain why cancellation request alone is insufficient.
\end{enumerate}

\section{Further reading}
\label{sec:ch74-reading}
The exception ABI explains personality routines and two-phase unwinding~\cite{itanium-2026-exceptions}.
Lua documents protected calls and error objects; Python documents its exception indicator and reference obligations~\cite{ierusalimschy-2026-lua,python-2026-capi}.
